% Chapter 6 — Beyond the Original Vision: Capabilities the Inventors Could Not Have Imagined
% Part of the OLDIA Thesis

\chapter{Beyond the Original Vision: Capabilities the Inventors Could Not Have Imagined}
\label{ch:inventors}

The researchers who built MYCIN, Hearsay, SOAR, STRIPS, GPS, CBR, DENDRAL, and ELIZA were among
the most capable scientists of their generation. Their ideas were correct. Their mathematics was
sound. Their theoretical contributions endure in the literature fifty years later. But they were
working on machines that ran at 1--5 MHz with 64 kilobytes to 8 megabytes of memory, and the
conceptual universe available to them was bounded by what was operationally possible on those
machines. This chapter examines what is now possible that was not only impossible then, but was
not in the conceptual universe of the original inventors at all.

The distinction matters. Some capabilities are simply faster versions of what the original
researchers envisioned. Others are qualitatively new --- they are not on the spectrum from slow to
fast but represent entirely new categories of epistemic activity that the original systems could
not have supported in any form. This chapter focuses on the latter.

\section{The Substrate Was Wrong by Six to Nine Orders of Magnitude}
\label{sec:substrate}

MYCIN ran on a PDP-10 at approximately 1--2 MHz with 64 KB of addressable memory. A single
MYCIN consultation --- loading the knowledge base, running the backward-chaining inference
engine, collecting certainty factors, and producing a recommendation --- took minutes. HEARSAY
ran on a Lisp machine with dedicated hardware at roughly the same speed. SOAR was developed on
VAX-11/780 systems running at approximately 5 MHz with 8 MB of memory; a SOAR problem-solving
episode involving subgoal resolution could take tens of seconds.

The wasm4pm MYCIN breed executes in 2--50 microseconds on 2026 hardware. This is not a
modest improvement. It is six to nine orders of magnitude. A million MYCIN consultations fit
inside a single second. The implications are not just quantitative. When an operation moves from
minutes to microseconds, it changes category. Operations that were practical once per session
become practical one hundred thousand times per second. That categorical shift opens epistemic
territory that the original researchers could not have mapped because they had no reason to
imagine it existed.

This is worth stating precisely: the ideas were correct. Shortliffe and Buchanan's certainty
factor algebra is correctly implemented in \texttt{production\_rules.rs}. Victor Lesser and Lee
Ercoli's KSA priority queue with opportunistic scheduling is correctly implemented in
\texttt{hearsay.rs}. Allen Newell's preference algebra is correctly implemented in
\texttt{soar.rs}. The implementations are not approximations or simplifications. They are
faithful to the original mathematical specifications. What is new is what becomes possible when
those correct ideas run on a substrate that is nine orders of magnitude more capable than the
one for which they were designed.

\section{Operational Falsifiability as a New Discipline}
\label{sec:falsifiability}

The original researchers used formal verification for algorithms. They proved correctness
properties about their systems in the mathematical sense: that STRIPS planning was complete under
certain conditions, that MYCIN's CF combination was monotone in particular configurations, that
SOAR's preference algebra was antisymmetric. These are classical theoretical computer science
contributions and they remain valid.

What they did not have, and could not have had, was any concept of \emph{runtime proof}.

Process mining as a discipline did not exist until Wil van der Aalst's work in the late 1990s.
The idea of deriving a structured event log from a running system and then applying conformance
checking algorithms to compare the observed behavior against a declared process model was not
published until the 2000s. OCEL 2.0 --- the Object-Centric Event Log standard that supports
multi-object process models --- was standardized in 2023.

The wasm4pm cognition layer implements operational falsifiability as a first-class architectural
constraint. Every breed execution passes through a choke point in \texttt{wasm.rs} that calls
\texttt{derive\_ocel(breed\_id, run\_id, \&[TraceStep])} and then
\texttt{validate\_ocel\_alignment(log, model)}. The \texttt{src/ocel.rs} file implements 632
lines of OCEL 2.0 infrastructure: \texttt{OcelEvent}, \texttt{OcelObject}, \texttt{OcelLog},
13 breed lifecycle models (\texttt{MYCIN\_MODEL}, \texttt{HEARSAY\_MODEL}, \texttt{CBR\_MODEL},
\texttt{GPS\_MODEL}, \texttt{STRIPS\_MODEL}, \texttt{PROLOG\_MODEL}, \texttt{SOAR\_MODEL},
\texttt{ELIZA\_MODEL}, \texttt{DENDRAL\_MODEL}, and four autoinstinct variants), a DFA phase
checker in \texttt{validate\_ocel\_alignment}, and a strictly monotone logical step enforcer in
\texttt{check\_temporal\_conformance}.

Non-conforming execution is refused. Not logged. Not warned about. Refused.

Shortliffe could not have conceived of this. Not because he lacked the intelligence, but because
the conceptual vocabulary did not exist. In 1976, proving that your algorithm was correct was
the state of the art. The idea that every individual run of a system would be checked against a
declared process model at runtime, and that the run would be rejected if it deviated, was not
in the research agenda of any laboratory. It represents a new discipline --- operational
falsifiability --- that the substrate of 1976 made not just impractical but inconceivable.

\section{Ensemble Reasoning at Operational Tempo}
\label{sec:ensemble}

At 2 microseconds per MYCIN execution, it is possible to run 500,000 MYCIN consultations per
second with systematic variations in certainty factor assignments. This is effectively a full
Bayesian sensitivity analysis over the CF algebra at interactive speed.

Consider what this enables. The original MYCIN ran one consultation at a time. A clinician
entered patient data, MYCIN fired its rules, and after minutes of processing, a recommendation
emerged with associated certainty factors. The certainty factors expressed the system's
confidence in its conclusion given the evidence. But they did not tell you how sensitive that
conclusion was to variations in the evidence weights. A clinician who wanted to know what would
happen if the certainty of a positive blood culture were reduced from 0.8 to 0.6 would have
to run a second consultation, wait minutes, and compare the outputs manually.

At 2 microseconds per execution, you can enumerate the entire grid of CF variations for a
consultation --- 100 values for each of 50 input CFs, producing $100^{50}$ points on the
sensitivity surface --- by sampling intelligently at rates that would produce publication-quality
sensitivity analyses in seconds. The landscape of the CF algebra becomes explorable rather than
a single point.

Hearsay is fundamentally an ensemble coordinator: its blackboard architecture was designed to
support multiple competing KSAs contributing partial solutions. Victor Lesser understood that
the value of the architecture was in the coordination of multiple hypotheses. But in 1977, the
KSAs that Hearsay coordinated were themselves expensive computations. The idea of running
13 different reasoning breeds simultaneously, treating each as a KSA contributing to a shared
blackboard, and selecting among their outputs using Hearsay's own priority scheduling, was
theoretically possible in Hearsay's architecture but operationally impossible.

The wasm4pm implementation makes this trivial. All 13 breeds share
\texttt{BreedInput}/\texttt{BreedOutput}. You can pipeline MYCIN's output as STRIPS's input,
pass STRIPS's plan to CBR for case-based refinement, and have HEARSAY coordinate the ensemble
--- in three lines. The original systems were islands of expertise separated by years of
integration work. Now composition is structural.

\section{Composition Without Integration Work}
\label{sec:composition}

The original systems were islands. MYCIN could not call STRIPS. SOAR could not query CBR.
Integrating two AI systems from different laboratories in the 1980s was a multi-year research
project involving custom data format translators, incompatible Lisp dialects, and fundamental
representational incompatibilities. The knowledge representation communities around frame-based
systems, rule-based systems, and planning systems were sometimes more incompatible with each
other than they were with the external world.

The wasm4pm architecture enforces a universal interface contract. Every breed accepts
\texttt{BreedInput} and returns \texttt{BreedOutput}. The \texttt{BreedInput} contract includes
a \texttt{breed} field (dispatching to the correct engine), a \texttt{contract} field containing
breed-specific inputs, and an optional \texttt{options.profile} for performance tuning. Rust's
\texttt{deny\_unknown\_fields} attribute on the input deserialization means that malformed inputs
are rejected at the boundary rather than producing silently incorrect behavior downstream.

This is not merely convenient. It is a structural guarantee that composition is always possible.
There is no integration work because there is no representation mismatch to bridge. MYCIN's
outputs can become STRIPS's inputs not because someone wrote an adapter, but because both
systems speak the same type language. This is what object-centric process modeling enables at
the representation level: a common object model makes inter-system composition a matter of
routing rather than translation.

The original researchers would have recognized this as the realization of a dream. The
Knowledge Interchange Format (KIF) project of the 1980s and 1990s spent a decade trying to
create a common language for AI systems to share knowledge. DARPA funded multiple initiatives
to achieve system interoperability. The problem was understood and considered important. What
was not available was the combination of a common type system, a WASM compilation target that
could run multiple systems in the same memory space, and the engineering discipline to enforce
a shared interface contract at compilation time.

\section{Process Mining Over Reasoning Chains}
\label{sec:process-mining}

Every breed execution appends to
\texttt{.wasm4pm/ocel/cognition/<breed>.jsonl}. After thousands of executions,
van der Aalst's process discovery algorithms can mine the actual process --- not the declared
process, but the one that actually ran --- from the accumulated event log.

The distinction between declared and actual process is the central insight of process mining.
A system's documentation describes what it is supposed to do. The event log records what it
actually did. Conformance checking quantifies the gap. The gap is always informative, often
surprising, and occasionally alarming.

For AI reasoning systems, this capability is genuinely new. The original researchers could
observe their systems' behavior by reading output traces. They could debug by examining rule
firings. They could measure performance by timing consultations. But they could not derive
a stochastic process model from a thousand executions and ask: what is the distribution over
execution paths? What variants occur? Are there hidden loops --- situations where the system
revisits states it has already processed? Is there concept drift --- evidence that the
system's behavior is changing over time, perhaps due to knowledge base modifications? Are
there execution patterns that suggest the certainty factor thresholds are miscalibrated?

Variant explosion detection, hidden loop discovery, concept drift monitoring: these are
not debugging techniques. They are forms of epistemic surveillance over the reasoning
process itself. The original researchers were building systems that would run occasionally
in laboratory settings. The idea of mining the distribution of reasoning patterns from
10 million MYCIN consultations was not conceivable because 10 million MYCIN consultations
at 1976 speeds would have taken approximately 317 years of continuous operation.

At 2 microseconds per execution, 10 million consultations take 20 seconds. The event
log that accumulates during routine operation becomes a scientific dataset about the
system's reasoning behavior. Process mining over that dataset is a new form of AI
introspection that did not exist and could not have existed in 1976.

\section{Cryptographic Provenance as Civil Infrastructure}
\label{sec:provenance}

The BLAKE3 receipt chain enables something the 1970s researchers would have found
incomprehensible: a globally unique identifier for a specific reasoning act that can be
presented as certified evidence in regulatory, legal, and contractual contexts.

MYCIN's 1976 consultation produced a printout. That printout recorded the system's
recommendation and supporting evidence. It was useful clinical documentation but it had no
formal properties. It could not be independently verified. It could not be reproduced from
stored inputs. It could not be checked for tampering. It was a document, not a proof.

The wasm4pm MYCIN produces a \texttt{run\_id} = BLAKE3(\textit{breed} $\|$ \textit{output\_hash})
that covers the entire reasoning process. The \texttt{output\_hash} =
BLAKE3(\texttt{serde\_json}(\texttt{BreedOutput})) covers the OCEL log embedded inside the
output. The \texttt{replay\_pointer} = \texttt{output\_hash[..16]} is a short reference that
can be presented in a compact form while the full receipt provides verifiable provenance.

What this means in practice: a MYCIN consultation performed by the wasm4pm implementation
produces a receipt that attests not just to the result but to the process. The receipt can
be stored. It can be presented to a regulator who asks: what was the system's reasoning at
this specific point in time? The answer is not a printout but a cryptographically attested
computation that can be independently reproduced from stored inputs and verified to produce
the identical receipt.

This is the difference between a doctor's handwritten note and a certified laboratory report
with a chain of custody. It is not merely a quality improvement. It is a categorical change
in what AI reasoning means for civil, regulatory, and legal purposes. The original researchers
were not building systems intended for regulated environments. They were building research
systems. The infrastructure for AI reasoning to operate in regulated environments --- with
auditable, reproducible, tamper-evident records --- did not exist and was not on the research
agenda.

\section{The Speed Enables Real-Time Conformance Checking}
\label{sec:realtime-conformance}

The operational falsifiability infrastructure described in Section~\ref{sec:falsifiability}
runs on every execution. This is only possible because the overhead of OCEL derivation and
conformance validation is negligible relative to the computation being validated.

In 1976, consider what runtime conformance checking would have required. Deriving a structured
event log from a MYCIN consultation would have required designing a log schema, instrumenting
the rule firing engine to emit log entries, writing the log to disk, and then running a
conformance checking algorithm over the log. Each of these steps --- instrumentation overhead,
disk I/O, log parsing, conformance algorithm execution --- would have consumed resources
comparable to or exceeding the original consultation. Checking whether the system was
conforming to its declared lifecycle model would have more than doubled the cost of running
the system.

In 2026, the \texttt{derive\_ocel} + \texttt{validate\_ocel\_alignment} pipeline runs in
microseconds. The overhead is noise against the measurement. This means conformance checking
is not an occasional audit but an every-execution invariant. It is structurally identical
to a type check: something that happens as a matter of course, whose failure is a hard error,
and whose cost is not budgeted separately because it is below the threshold of concern.

This is a new relationship between a system and its specification. The system does not merely
claim to conform to its declared lifecycle model. It proves conformance on every execution or
refuses to complete. The specification is not a document that could become stale. It is an
executable constraint enforced at the choke point.

\section{The Unfakeable Test Suite as a New Kind of Specification}
\label{sec:specification}

The 655-test suite --- 480 Rust tests passing, 175 TypeScript tests passing, zero failures ---
is not merely quality assurance infrastructure. It is a formal specification of what each
algorithm must do, expressed in a form that falsifies incorrect implementations at the level
of mathematical theorems.

Consider the Prolog grandparent test. The query
\texttt{grandparent(?0,?2) :- parent(?0,?1), parent(?1,?2)} must derive that carol has a
grandparent from the facts parent(alice, bob) and parent(bob, carol). This test is not
checking that the system produces a particular output for a particular input. It is checking
that Robinson flat-term unification with shared variable bindings is correctly implemented.
If \texttt{?1} is not properly shared between the two parent clauses --- if the implementation
binds \texttt{?1} in the first clause independently from \texttt{?1} in the second --- the
grandparent inference fails. The test is true if and only if the mathematical specification
of SLD resolution with shared variables is correctly implemented. No other implementation
produces the correct output.

This property --- that certain tests are true if and only if the mathematical specification
is correctly implemented --- is what makes them unfakeable. They are not integration tests
that check whether the system produces acceptable outputs. They are falsification tests that
check whether specific mathematical properties hold. The MYCIN CF strict boundary test
(CF = 0.2 rejected, CF = 0.201 accepted) is unfakeable without correctly implementing
the Shortliffe-Buchanan combination formula with a strict inequality threshold. The CBR
Jaccard identity test (identical query and case must produce similarity = 1.0 exactly) is
unfakeable without correctly implementing Jaccard similarity computation. The SOAR
antisymmetry test (circular worse-than preferences must produce deterministic termination)
is unfakeable without implementing the preference algebra with the antisymmetry property
that Newell specified.

The original researchers wrote papers. The papers specified the algorithms in mathematical
notation and English prose. The specifications were correct. But they were not executable.
A reader of Shortliffe and Buchanan (1975) who wanted to verify that their implementation
of MYCIN's CF algebra was correct had to derive test cases from the mathematical formulas
by hand, run the system on those test cases, and compare the outputs manually. The process
was laborious, error-prone, and incomplete.

A test suite that is a formal specification is a new kind of artifact. It is not documentation
of behavior. It is a machine-executable oracle that either confirms correct implementation
or rejects it. For the wasm4pm cognition layer, the test suite is the specification. The
papers by Shortliffe, Newell, Lesser, Fikes, Hart, Nilsson, and Aamodt provide the
mathematical foundations. The test suite is the operational embodiment of those foundations ---
the form in which they can be applied to arbitrary implementations to determine correctness.

The original researchers would have recognized the value immediately. Shortliffe spent years
defending MYCIN's CF algebra against critics who questioned whether it was correctly
implemented in clinical validation studies. An unfakeable test suite would have settled those
arguments definitively. What was not available in 1976 was the combination of fast execution
(making exhaustive test suites practical), modern testing infrastructure (making property-based
tests and adversarial oracles straightforward to write), and the accumulated understanding of
what mathematical properties distinguish correct from incorrect implementations.

\section{Summary: A New Epistemic Category}
\label{sec:summary}

The capabilities described in this chapter share a common structure: they are not quantitative
improvements on what the original researchers envisioned. They are qualitatively new epistemic
activities that became possible when correct ideas from 1976--1987 ran on a substrate nine
orders of magnitude more capable than the one for which they were designed.

Operational falsifiability: the original researchers verified algorithms. We verify every
execution against a declared process model.

Ensemble reasoning at operational tempo: the original researchers understood ensemble
architectures. We can execute them with 500,000 iterations per second, making the full
sensitivity surface of the CF algebra explorable interactively.

Composition without integration work: the original researchers dreamed of AI system
interoperability. We have it structurally, enforced at compilation time.

Process mining over reasoning chains: the original researchers observed system behavior
by reading traces. We mine stochastic process models from millions of executions and
apply conformance checking, variant analysis, and concept drift detection.

Cryptographic provenance: the original researchers produced printouts. We produce
cryptographically attested computation records with reproducible receipts and chain-of-custody
properties suitable for regulated environments.

The unfakeable test suite as specification: the original researchers wrote papers. We write
falsifying oracles that are true if and only if the mathematical specification is correctly
implemented.

These are not improvements to the original systems. They are new capabilities that the
original systems' architectures did not support and the original researchers did not envision,
because the conceptual vocabulary for most of them did not exist in 1976 and the operational
cost of all of them was prohibitive until recently.

The inventors were right about the ideas. They could not have imagined what would become
possible when those ideas ran on a substrate nine orders of magnitude beyond what they had.
